%%%%%%%%%%%%%%%%%%%%%%%%%%%%%%%%%%%%%%%%%%%%%%%%%%%%%%%%%%%%%%%%%%%%%%%%%%%%%%%%
%%  PRÄAMBEL – gemeinsam für Endabgabe und alle Zwischenabgaben.
%%  Teil 1 ist 1:1 die Präambel der offiziellen Gerüste (nicht verändern),
%%  Teil 2 ergänzt nur Pakete für Gleichungen, Bilder, PDF-Einbindung usw.,
%%  Teil 3 enthält die Hilfsmakros dieser Vorlage.
%%  Die \documentclass-Zeile steht (unverändert offiziell) in jeder Hauptdatei.
%%%%%%%%%%%%%%%%%%%%%%%%%%%%%%%%%%%%%%%%%%%%%%%%%%%%%%%%%%%%%%%%%%%%%%%%%%%%%%%%

%%%%%%%%%%%%%%%%%%%%%%%%%%%%% Teil 1: offizielle Präambel %%%%%%%%%%%%%%%%%%%%%%
\usepackage[verbose,a4paper,bottom=25mm,top=25mm]{geometry} %showframe,
\usepackage{scrlayer-scrpage}
\usepackage[T1]{fontenc}
\usepackage[utf8]{inputenc}
\usepackage{ngerman}
\usepackage{color}
\usepackage{amsfonts}
\usepackage{amssymb}
\usepackage{amsthm}
%%%%%%%%%%%%%%%%%%%%%%%%% Biblatex (offizielle Optionen, Backend Biber) %%%%%%%%%%
%% Overleaf erkennt biber automatisch. Zitierstil alphabetic: [Sch02]
\usepackage[sorting=nyt, backend=biber, citestyle=alphabetic, style=alphabetic,
            autocite=inline, labelalpha=true, hyperref=auto, backref=true,
            language=ngerman]{biblatex}
\addbibresource{literatur.bib}

%%%%%%%%%%%%%%%%%%%%%%%%%%%%% Teil 2: Zusatzpakete %%%%%%%%%%%%%%%%%%%%%%%%%%%%%%
\usepackage{amsmath}               % \eqref, align, nummerierte Gleichungen
\usepackage{xcolor}
\definecolor{darkgreen}{rgb}{0,0.5,0}   % offizielle Hinweisfarbe
\usepackage{graphicx}              % \includegraphics
\graphicspath{{abbildungen/}}
\usepackage{booktabs}              % schöne Tabellen
\usepackage[locale=DE]{siunitx}    % \SI{3.4}{\kilo\hertz}, \num{1.2e3}
\usepackage{pdfpages}              % \includepdf für Anhänge (Paper, Laborjournal, Folien)
\usepackage{pdflscape}             % Querformat: \begin{landscape}...\end{landscape}
\usepackage{pgfgantt}              % Gantt-Diagramm für den Arbeitsplan
\usepackage{environ}
\usepackage[savepos,abspage]{zref}
\usepackage[hidelinks]{hyperref}   % klickbare Verweise; zuletzt laden

%% Zitierregel (Marti): bei Journalartikeln DOI angeben, aber keine Webadresse.
%% -> url/urldate werden bei @article automatisch unterdrückt.
\AtEveryBibitem{%
  \ifentrytype{article}{\clearfield{url}\clearfield{urlyear}%
    \clearfield{urlmonth}\clearfield{urlday}}{}}

%% Kopf- und Fußzeile wie offizielles Gerüst (\ohead wird in der Hauptdatei gesetzt)
\pagestyle{scrheadings}
\ihead{\thepage}
\ofoot{\pagemark}
\ifoot{\zubearbeiten{\IhrName}}

%%%%%%%%%%%%%%%%%%%%%%%%%%%%% Teil 3: Hilfsmakros %%%%%%%%%%%%%%%%%%%%%%%%%%%%%%%
\makeatletter

%% --- offizieller Farbschalter (Zeile 12-15 der Gerüste), automatisch per Modus
\ifAbgabe
  \Hinweisefalse
  \newcommand{\zubearbeiten}[1]{\textcolor{black}{#1}}
\else
  \newcommand{\zubearbeiten}[1]{\textcolor{darkgreen}{#1}}
\fi

%% --- Grüne Hilfstexte: verschwinden im Abgabemodus vollständig
%%     \begin{hinweise} ... \end{hinweise}   (darf itemize enthalten)
%%     \hinweis{kurzer Text}
\NewEnviron{hinweise}{\ifHinweise\par{\color{darkgreen}\small\BODY\par}\fi}
\newcommand{\hinweis}[1]{\ifHinweise\par{\color{darkgreen}\small #1\par}\fi}

%% --- Platzhalter für eigenen Text: bleibt IMMER rot sichtbar (auch im
%%     Abgabemodus), damit nichts Unfertiges abgegeben wird.
\newcounter{platzhalter}
\DeclareRobustCommand{\platzhalter}[1]{%
  \stepcounter{platzhalter}%
  \textcolor{red}{\textbf{[}\,#1\,\textbf{]}}%
  \ifAbgabe\PackageWarning{Portfolio}{Platzhalter noch vorhanden: #1}\fi}
\AtEndDocument{%
  \ifnum\value{platzhalter}>0
    \PackageWarningNoLine{Portfolio}{Es sind noch \theplatzhalter\space
      Platzhalter im Dokument!}%
  \fi
  \immediate\write\@auxout{\string\gdef\string\pf@platzhalteralt{\theplatzhalter}}}
\providecommand{\pf@platzhalteralt}{?}

%% --- Umfangsmessung
%%   \begin{umfang}{Bezeichnung}{min}{max} ... \end{umfang}
%%   misst die Länge des eingeschlossenen Textes in Seiten (über die
%%   Position auf der Seite, auf 0.05 Seiten genau) und zeigt sie im Entwurf
%%   in einem Kasten an. Achtung: Abbildungen/Tabellen zählen nach den Regeln
%%   NICHT zur Länge, werden hier aber mitgemessen -> dann abziehen.
\zref@addprop{savepos}{abspage}
\newcounter{pf@umfang}
\newcommand*{\pf@umfangliste}{}
\newcommand*{\pf@umfangeintrag}[4]{% Name, min, max, ist
  \g@addto@macro\pf@umfangliste{\pf@umfangzeile{#1}{#2}{#3}{#4}}}
\newcommand*{\pf@setfarbe}[3]{% min max ist -> \pf@farbe
  \def\pf@farbe{darkgreen}%
  \ifdim #3pt<#1pt \def\pf@farbe{red}\fi
  \ifdim #3pt>#2pt \def\pf@farbe{red}\fi}
\newcommand*{\pf@umfangzeile}[4]{%
  #1 & #2 -- #3 & \pf@setfarbe{#2}{#3}{#4}\textcolor{\pf@farbe}{\textbf{#4}}\\}
\newenvironment{umfang}[3]{%
  \stepcounter{pf@umfang}%
  \edef\pf@uid{pfumf\the\value{pf@umfang}}%
  \def\pf@uname{#1}\def\pf@umin{#2}\def\pf@umax{#3}%
  \par\zsavepos{\pf@uid a}\nobreak
}{%
  \par\zsavepos{\pf@uid e}%
  \ifAbgabe\else
    \zref@refused{\pf@uid a}\zref@refused{\pf@uid e}%
    \edef\pf@ist{\fpeval{round(
        (\zref@extractdefault{\pf@uid e}{abspage}{0}
        -\zref@extractdefault{\pf@uid a}{abspage}{0})
        +(\zref@extractdefault{\pf@uid a}{posy}{0}
        -\zref@extractdefault{\pf@uid e}{posy}{0})/\number\textheight ,2)}}%
    \protected@write\@auxout{}{\string\pf@umfangeintrag{\pf@uname}{\pf@umin}{\pf@umax}{\pf@ist}}%
    \par\smallskip\noindent
    \pf@setfarbe{\pf@umin}{\pf@umax}{\pf@ist}%
    \fcolorbox{\pf@farbe}{white}{\scriptsize
      Umfang \emph{\pf@uname}: $\approx$\,\pf@ist~Seiten
      (Soll: \pf@umin\,--\,\pf@umax~Seiten, ohne Titel/Abb./Tab./Legenden)}%
    \par\smallskip
  \fi}

%% --- Kontrollseite (nur im Entwurf; ersetzt das grüne Kapitel
%%     "Hinweise zur Arbeit mit dieser Datei", das abgegeben nicht vorkommen darf)
\newcommand{\kontrollseite}{%
  \ifAbgabe\else
  \chapter*{Kontrollseite (erscheint nur im Entwurf)}
  {\color{darkgreen}
  Modus: \textbf{Entwurf}. Für die Abgabe in \texttt{einstellungen.tex}
  \texttt{\textbackslash Abgabetrue} setzen. Dann wird alles Grüne schwarz bzw.\
  entfernt (Regel: \enquote{Abgaben mit grüner Farbe oder stehen gebliebenen
  Hilfen ergeben einen Abzug}).\par\medskip
  \textbf{Offene Platzhalter} (letzter Durchlauf): \pf@platzhalteralt\par\medskip
  \textbf{Umfangskontrolle} (Werte aus dem letzten Kompilierdurchlauf;
  Hilfstexte stehen außerhalb der Messung):\par\smallskip
  \ifx\pf@umfangliste\@empty
    \emph{noch keine Messwerte -- nochmals kompilieren.}
  \else
    {\small\begin{tabular}{p{0.55\textwidth}cc}
      \toprule Teil & Soll [Seiten] & Ist [Seiten]\\\midrule
      \pf@umfangliste
      \bottomrule
    \end{tabular}}
  \fi\par\medskip
  %% Checkliste – erscheint nur auf der Kontrollseite im Entwurf.
%% Quellen: Portfolios im Projektpraktikum 2026-2027 (Marti, 09.07.2026),
%% offizielle Gerüste 2026-2027, Zitierregeln (31.08.2026),
%% Abbildungen und Gleichungen (Limmer 2007), Moodle-Terminliste.

\textbf{Fristen (Moodle; nächste Abgabe wird erst nach der vorherigen freigeschaltet!)}\par\smallskip
{\small\begin{tabular}{@{}ll@{}}
21.\,10.\,2026 & Motivationsformular (ohne Reflexion) -- obligatorisch\\
28.\,10.\,2026 & Arbeitsplan + Reflexion -- obligatorisch\\
11.\,11.\,2026 & Wissenschaftliche Arbeit + Reflexion -- obligatorisch\\
09.\,12.\,2026 & Protokoll eines Praktikumstages + Reflexion -- obligatorisch\\
20.\,01.\,2027 & Messergebnis + Reflexion -- obligatorisch\\
23.\,01.\,2027 & Späteste Anmeldung mündliche Leistung (Portal)\\
27.\,01.\,2027 & Folien des Vortrags + Materialien für Mitstudierende\\
28.\,01.\,2027 & Vorträge (Prüfung, Anwesenheitspflicht)\\
07.\,02.\,2027 & Folien mit Kommentar + Reflexion (vorgeschlagen, freiwillig)\\
21.\,02.\,2027 & Aussage aus anderem Vortrag + Reflexion (vorgeschlagen, freiwillig)\\
10.\,03.\,2027 & Späteste Anmeldung schriftliche Leistung (Portal)\\
\textbf{15.\,03.\,2027} & \textbf{Endabgabe Portfolio -- verspätet = nicht bestanden}\\
\end{tabular}}\par\smallskip
{\small Praktikumstage (Google Kalender): donnerstags 10--12\,Uhr, O26/233 bzw.\ O26/234 --
15.\,10.\ bis 17.\,12.\,2026 und 14.\,01.\ bis 11.\,02.\,2027; keine Termine 18.\,12.--13.\,01.
Die Anmeldetermine im Hochschulportal (23.\,01./10.\,03.) stammen aus dem Plan von Petros -- im Portal prüfen!}
\par\medskip

\textbf{Vor jeder Abgabe prüfen}
\begin{itemize}\itemsep0pt
\item \texttt{\textbackslash Abgabetrue} gesetzt, keine grüne Farbe, keine Hilfen, \emph{0 Platzhalter}.
\item Deckblatt vollständig: Name, Vorname, Matrikelnummer, Studiengang, E-Mail, Thema (Projekt + Schwerpunkt).
\item Inhaltsverzeichnis mit Seitenzahlen, auch der Anhänge; durchgehend Seitenzahlen.
\item Längen eingehalten (siehe Tabelle oben; Abbildungen/Tabellen/Legenden zählen nicht mit):
      Einleitung 2--4, Beschreibung $\tfrac14$--$\tfrac12$, Diskussion der Physik 1--3,
      Reflexion 1--2, Gesamtreflexion 3--6 Seiten. Schriftgröße 12\,pt (fest eingestellt).
\item Beschreibung und Reflexion sind \emph{keine Nacherzählung}; die Leitfragen werden in Fließtext
      beantwortet (nicht in Einzelsätzen), erkennbar zugeordnet.
\item Paper (Einlage 2) wird in den späteren Einlagen zur Argumentation verwendet; Gleichungen werden
      zur Argumentation genutzt.
\item Peer-Feedback (Kommilitonen) eingeholt, aufgegriffen -- und zitiert. Petros' Arbeitsplan als Vorlage
      zitieren (Key \texttt{Petros2026Arbeitsplan}); Schwerpunkt-Abgrenzung mit Janosch festhalten.
\item Jede Übernahme zitiert: Literatur, Bilder (\texttt{\textbackslash cite[Abb.~3.4]\{key\}}),
      Gleichungen (\texttt{\textbackslash cite[Gl.~(3)]\{key\}}), Aussagen von Personen, Anleitungen,
      Datenblätter, Webseiten (mit Abrufdatum), \textbf{jede KI-Nutzung mit System und ALLEN Prompts}.
\item Journalartikel: Autoren, Titel, Journal, Band, Seiten, Jahr, DOI (nicht als Link).
\item Abbildungen: Bildunterschrift, im Text referenziert (\texttt{\textbackslash ref}), Schrift
      $\geq 10$\,pt, Achsen mit Einheit als \enquote{$f$ / kHz} (nicht \enquote{[kHz]}), ggf.\ Residuenplot.
\item Tabellen: Überschrift, im Text referenziert; Messdaten nicht als lange Listen im Text, sondern im Anhang.
\item Wichtige Gleichungen nummeriert, alle Symbole definiert und einheitlich verwendet.
\item Dateiname für Moodle: \texttt{Manevski\_WelcheEinlage.pdf} (nur PDF oder DOCX).
\item Bis auf den Arbeitsplan müssen sich alle Einlagen von denen der Partnerin/des Partners unterscheiden.
\end{itemize}

  }%
  \fi}
\providecommand{\enquote}[1]{\glqq #1\grqq}

%% --- Anhang aus PDF mit durchgehenden Seitenzahlen (Regel: "Die Arbeit ist
%%     durchgängig mit Seitenzahlen versehen ... auch der Anhänge")
%%   \anhangpdf{Kapiteltitel}{label}{datei.pdf}           alle Seiten
%%   \anhangpdf[pages={1-3}]{Kapiteltitel}{label}{datei.pdf}
%%   Die Kapitelüberschrift steht auf einer eigenen Seite, danach folgen die
%%   PDF-Seiten (auf 82 % verkleinert, damit Kopf-/Fußzeile frei bleiben).
\newcommand{\anhangpdf}[4][pages=-]{%
  \chapter{#2}\label{#3}%
  \includepdf[#1,scale=0.82,pagecommand={\thispagestyle{scrheadings}}]{#4}}
%% Gleiches ohne neues Kapitel (für weitere PDFs im selben Anhang)
\newcommand{\pdfseiten}[2][pages=-]{%
  \includepdf[#1,scale=0.82,pagecommand={\thispagestyle{scrheadings}}]{#2}}

\makeatother

%% Antwortfeld im Motivationsformular (Anhang der Endabgabe)
\newcommand{\antwort}[1]{\par\smallskip #1\par\bigskip}
